\documentclass[12pt]{article}
\usepackage[utf8]{inputenc}
\usepackage[T1]{fontenc}
\usepackage{amsmath, amssymb, geometry, booktabs}
\usepackage{enumitem}
\geometry{margin=1in}
\setlength{\parindent}{0pt}
\renewcommand{\arraystretch}{1.2}

\title{Solutions -- Final Exam Review}
\author{ECON 308 -- International Economic Relations}
\date{Fall 2025}

\begin{document}
\maketitle

% ============================================================
\section*{Problem 1 — Ricardian Model}

\textbf{Given:} Two countries, Home and Foreign, produce wheat (W) and cloth (C).  
Unit labor requirements:

\begin{center}
\begin{tabular}{lcc}
\toprule
Country & $a_{LW}$ & $a_{LC}$ \\
\midrule
Home    & 4         & 2        \\
Foreign & 5         & 1        \\
\bottomrule
\end{tabular}
\end{center}

Home labor: $L_H = 200$; Foreign labor: $L_F = 150$.

\begin{enumerate}[label=(\arabic*)]

\item \textbf{Autarky relative prices $P_W/P_C$.}

In the Ricardian model with one factor (labor), the opportunity cost of wheat in terms of cloth equals the ratio of unit labor requirements:
\[
\left(\frac{P_W}{P_C}\right)^{\text{aut}} = \frac{a_{LW}}{a_{LC}}.
\]

For Home:
\[
\left(\frac{P_W}{P_C}\right)_H^{\text{aut}} = \frac{4}{2} = 2.
\]

For Foreign:
\[
\left(\frac{P_W}{P_C}\right)_F^{\text{aut}} = \frac{5}{1} = 5.
\]

\[
\boxed{\left(P_W/P_C\right)_H^{\text{aut}} = 2,\quad \left(P_W/P_C\right)_F^{\text{aut}} = 5.}
\]

\item \textbf{Comparative advantage.}

Opportunity cost of 1 unit of wheat in terms of cloth:

Home:
\[
OC_H(W) = \frac{a_{LW}}{a_{LC}} = \frac{4}{2} = 2 \text{ units of cloth per unit of wheat.}
\]

Foreign:
\[
OC_F(W) = \frac{a_{LW}}{a_{LC}} = \frac{5}{1} = 5.
\]

Lower opportunity cost in wheat: Home (2 $<$ 5) $\Rightarrow$ Home has \textbf{comparative advantage in wheat}.  
Similarly, opportunity cost of cloth is the inverse:

Home:
\[
OC_H(C) = \frac{a_{LC}}{a_{LW}} = \frac{2}{4} = 0.5 \text{ wheat per cloth.}
\]
Foreign:
\[
OC_F(C) = \frac{1}{5} = 0.2 \text{ wheat per cloth.}
\]

Lower opportunity cost in cloth: Foreign $\Rightarrow$ \textbf{Foreign has comparative advantage in cloth.}

\[
\boxed{\text{Home: CA in wheat; Foreign: CA in cloth.}}
\]

\item \textbf{World relative price $P_W/P_C = 1.5$: exports and production bundles.}

Formally, firms choose the good that maximizes revenue per unit of labor:
\[
\text{Revenue per labor in } W: \frac{P_W}{a_{LW}},\quad
\text{in } C: \frac{P_C}{a_{LC}}.
\]

Normalize $P_C = 1$, so $P_W = 1.5$.

\emph{Home:}
\[
\frac{P_W}{a_{LW}} = \frac{1.5}{4} = 0.375, \quad 
\frac{P_C}{a_{LC}} = \frac{1}{2} = 0.5.
\]
So cloth gives higher revenue per unit of labor. Home would fully specialize in cloth at this relative price.

\emph{Foreign:}
\[
\frac{P_W}{a_{LW}} = \frac{1.5}{5} = 0.3, \quad 
\frac{P_C}{a_{LC}} = \frac{1}{1} = 1.
\]
Again cloth yields higher returns; Foreign also specializes in cloth.

Thus, at $P_W/P_C = 1.5$:
\[
\boxed{\text{Both Home \emph{and} Foreign would specialize in cloth.}}
\]

This relative price is actually \emph{inconsistent} with trade equilibrium in the basic Ricardian model, because no one wants to produce wheat at such a low wheat price. In a consistent free-trade equilibrium, $P_W/P_C$ must satisfy:
\[
\frac{a_{LW}}{a_{LC}}\Big|_H < \frac{P_W}{P_C} < \frac{a_{LW}}{a_{LC}}\Big|_F
\Rightarrow 2 < \frac{P_W}{P_C} < 5.
\]
For example, if $P_W/P_C=3$, Home would specialize in wheat and Foreign in cloth, and trade would be feasible. The question’s price (1.5) is useful pedagogically to illustrate this inconsistency.

\item \textbf{Gains from trade intuition.}

Even if one country were absolutely less productive in both goods, specialization according to \emph{comparative} advantage allows:
\begin{itemize}
\item Each country to specialize where its \emph{relative} productivity is highest.
\item Total world output to rise for at least one good.
\item Each country, by trading at world prices between its own autarky relative price and its partner’s, to reach a higher consumption possibility set and indifference curve.
\end{itemize}

So in a valid free-trade equilibrium (with $2 < P_W/P_C < 5$), both Home and Foreign can consume bundles that are infeasible under autarky.

\end{enumerate}

% ============================================================
\section*{Problem 2 — Specific Factors Model}

\textbf{Given:}  
Home produces manufactures (M) and food (F). Labor $L$ is mobile; capital $K$ is specific to M; land $T$ is specific to F.

\[
Q_M = 5K^{1/2}L_M^{1/2}, \qquad Q_F = 4T^{1/2}L_F^{1/2},
\]
with $L_M + L_F = L = 200$, $K = 100$, $T = 100$.

\begin{enumerate}[label=(\arabic*)]

\item \textbf{Marginal product of labor in each sector.}

For M:
\[
Q_M = 5K^{1/2}L_M^{1/2}.
\]
Differentiate w.r.t.\ $L_M$:
\[
MPL_M = \frac{\partial Q_M}{\partial L_M}
       = 5K^{1/2} \cdot \frac{1}{2}L_M^{-1/2}
       = \frac{5}{2}K^{1/2}L_M^{-1/2}.
\]

For F:
\[
Q_F = 4T^{1/2}L_F^{1/2}
\Rightarrow MPL_F = \frac{\partial Q_F}{\partial L_F}
       = 4T^{1/2}\cdot\frac{1}{2}L_F^{-1/2}
       = 2T^{1/2}L_F^{-1/2}.
\]

\[
\boxed{MPL_M = \frac{5}{2}K^{1/2}L_M^{-1/2},\quad MPL_F = 2T^{1/2}L_F^{-1/2}.}
\]

\item \textbf{Equilibrium condition for labor allocation.}

In the specific factors model, with competitive firms:
\[
w = P_M \cdot MPL_M = P_F \cdot MPL_F,
\]
and $L_M + L_F = L$.

Thus the equilibrium allocation $(L_M,L_F)$ solves:
\[
P_M \cdot \frac{5}{2}K^{1/2}L_M^{-1/2}
=
P_F \cdot 2T^{1/2}L_F^{-1/2}
\quad \text{subject to} \quad L_F = L - L_M.
\]

This gives the interior allocation where workers are indifferent between sectors.

\item \textbf{Effect of a rise in the relative price of manufactures, $P_M/P_F$.}

Suppose $P_M$ rises while $P_F$ is fixed.

\begin{itemize}
\item The equality $P_M MPL_M = P_F MPL_F$ implies that, for given $L_M,L_F$, the left-hand side rises, so the wage $w$ initially rises.
\item As $w$ rises, labor moves into manufacturing (M), raising $L_M$ and lowering $L_F$, so:
\[
MPL_M \downarrow \quad (\text{due to diminishing returns in M}),
\quad
MPL_F \uparrow \quad (\text{less labor in F}).
\]
\item The new equilibrium has:
\[
w = P_M MPL_M \quad \text{(higher $P_M$, lower MPL$_M$)}, \\
w = P_F MPL_F \quad \text{(fixed $P_F$, higher MPL$_F$)}.
\]
Result:
\begin{itemize}
\item \textbf{Nominal wage} $w$ rises, but by less than $P_M$.
\item Hence \textbf{real wage in terms of M} ($w/P_M$) falls, while \textbf{real wage in terms of F} ($w/P_F$) rises.
\end{itemize}
\end{itemize}

\item \textbf{Who gains, who loses?}

\begin{itemize}
\item Capital (specific to M):  
  $r_K = P_M \cdot MPK_M$ rises more than the wage; M expands, and the specific factor in the booming sector gains \textbf{unambiguously in real terms}.
\item Land (specific to F):  
  $r_T = P_F \cdot MPT_F$ \textbf{falls in real terms}. The specific factor in the contracting sector loses.
\item Labor (mobile factor):  
  Ambiguous: real wage rises in terms of F but falls in terms of M. The net effect depends on each worker's consumption bundle.
\end{itemize}

\[
\boxed{\text{Specific factor in the rising-price sector gains;}}
\]
\[
\boxed{\text{Specific factor in the other sector loses; labor is ambiguous.}}
\]
\end{enumerate}

% ============================================================
\section*{Problem 3 — Heckscher--Ohlin Model}

\textbf{Given:} Two goods, machinery (M, capital-intensive) and textiles (T, labor-intensive).  
Factor endowments:

\begin{center}
\begin{tabular}{lcc}
\toprule
Country & $K$ & $L$ \\
\midrule
Home    & 900 & 400 \\
Foreign & 300 & 600 \\
\bottomrule
\end{tabular}
\end{center}

\begin{enumerate}[label=(\arabic*)]

\item \textbf{Factor abundance.}

Compute $K/L$:

Home:
\[
(K/L)_H = 900/400 = 2.25.
\]
Foreign:
\[
(K/L)_F = 300/600 = 0.5.
\]

Home has higher $K/L$, so is \textbf{capital-abundant}.  
Foreign has lower $K/L$, so is \textbf{labor-abundant}.

\[
\boxed{\text{Home capital-abundant, Foreign labor-abundant.}}
\]

\item \textbf{Trade pattern (H--O theorem).}

Heckscher–Ohlin: a country exports the good that uses its \emph{abundant} factor intensively.

\begin{itemize}
\item Home (capital-abundant) exports the \textbf{capital-intensive good}, machinery (M).
\item Foreign (labor-abundant) exports the \textbf{labor-intensive good}, textiles (T).
\end{itemize}

\[
\boxed{\text{Home exports M, imports T; Foreign exports T, imports M.}}
\]

\item \textbf{Winners and losers (Stolper–Samuelson).}

When a country opens to trade and the relative price of its export good rises:

\begin{itemize}
\item In Home, $P_M/P_T$ rises (M is exported).  
  Stolper–Samuelson $\Rightarrow$ the real return to the factor used intensively in M (capital) rises; the real wage (labor) falls.
  \[
  \boxed{\text{In Home: capital owners gain; workers lose.}}
  \]
\item In Foreign, $P_T/P_M$ rises (T is exported), so the real wage (labor) rises and the real return to capital falls.
  \[
  \boxed{\text{In Foreign: workers gain; capital owners lose.}}
  \]
\end{itemize}

\item \textbf{Rybczynski effect of higher $L$ in Home.}

Rybczynski theorem: at fixed goods prices, an increase in one factor endowment increases output of the good that uses that factor intensively and reduces output of the other good.

In Home:

\begin{itemize}
\item If $L$ increases (labor endowment rises), while $K$ and prices are fixed, output of the labor-intensive good (T) rises, and output of the capital-intensive good (M) falls.
\item The economy ``tilts'' toward more textiles for given goods prices.
\end{itemize}

\[
\boxed{\text{In Home: higher $L$} \Rightarrow Q_T \uparrow,\; Q_M \downarrow \;\text{(at given prices).}}
\]

\end{enumerate}

% ============================================================
\section*{Problem 4 — Standard Trade Model}

\textbf{Given:} PPF:
\[
A^2 + 4E^2 = 1600,
\]
where $A$ = apples, $E$ = electronics.

World relative price:
\[
P_E/P_A = 2.
\]

Normalize $P_A=1$ so $P_E=2$. Preferences at the consumption point satisfy $E = 0.25A$.

\begin{enumerate}[label=(\arabic*)]

\item \textbf{Production point maximizing value.}

Maximize value:
\[
V = P_A A + P_E E = A + 2E
\]
subject to PPF: $A^2 + 4E^2 = 1600$.

The slope of the PPF:
\[
A^2 + 4E^2 = 1600 \Rightarrow 2A\,dA + 8E\,dE = 0 \Rightarrow \frac{dE}{dA} = -\frac{A}{4E}.
\]
Slope of the isovalue line: $-\frac{P_A}{P_E} = -\frac{1}{2}$.

Tangency (interior solution):
\[
-\frac{A}{4E} = -\frac{1}{2} \Rightarrow \frac{A}{4E} = \frac{1}{2} \Rightarrow A = 2E.
\]

Substitute into PPF:
\[
(2E)^2 + 4E^2 = 4E^2 + 4E^2 = 8E^2 = 1600 \Rightarrow E^2 = 200 \Rightarrow E_P = \sqrt{200} = 10\sqrt{2}.
\]
Then
\[
A_P = 2E_P = 20\sqrt{2}.
\]

\[
\boxed{(A_P,E_P) = (20\sqrt{2},\,10\sqrt{2}).}
\]

\item \textbf{Consumption bundle.}

Income (value of production):
\[
V = A_P + 2E_P = 20\sqrt{2} + 2\cdot 10\sqrt{2} = 40\sqrt{2}.
\]

Budget line at world prices:
\[
A + 2E = V = 40\sqrt{2}.
\]

Preferences at the chosen bundle satisfy $E = 0.25A$. Substitute:
\[
A + 2(0.25A) = A + 0.5A = 1.5A = 40\sqrt{2}.
\]
Hence
\[
A_C = \frac{40}{1.5}\sqrt{2} = \frac{80}{3}\sqrt{2}, \quad
E_C = 0.25A_C = \frac{20}{3}\sqrt{2}.
\]

\[
\boxed{(A_C,E_C) = \left(\frac{80}{3}\sqrt{2},\,\frac{20}{3}\sqrt{2}\right).}
\]

\item \textbf{Exports and imports.}

Compare production and consumption:

\[
E\text{ exports}: X_E = E_P - E_C = 10\sqrt{2} - \frac{20}{3}\sqrt{2}
= \frac{30 - 20}{3}\sqrt{2} = \frac{10}{3}\sqrt{2} >0.
\]

\[
A\text{ imports}: M_A = A_C - A_P = \frac{80}{3}\sqrt{2} - 20\sqrt{2}
= \frac{80 - 60}{3}\sqrt{2} = \frac{20}{3}\sqrt{2} >0.
\]

\[
\boxed{\text{Country exports electronics and imports apples.}}
\]

Check trade balance at world prices ($P_A=1,P_E=2$):

Export value:
\[
P_E X_E = 2 \cdot \frac{10}{3}\sqrt{2} = \frac{20}{3}\sqrt{2}.
\]
Import value:
\[
P_A M_A = 1 \cdot \frac{20}{3}\sqrt{2} = \frac{20}{3}\sqrt{2}.
\]
Balanced trade holds.

\item \textbf{Effect of a rise in $P_E/P_A$.}

If $P_E/P_A$ rises (electronics become relatively more valuable):
\begin{itemize}
\item \textbf{Production:} The isovalue line becomes steeper; the new tangency with the PPF shifts towards more $E$ and fewer $A$. Export sector expands.
\item \textbf{Consumption:} The trade (budget) line becomes steeper and passes through the new production point, allowing the country to reach a higher indifference curve (since it is an exporter of $E$ and its terms of trade improve).
\item \textbf{Welfare:} With no distortions, an improvement in the terms of trade (price of export good rises) increases national welfare: the country can trade the same export volume for more imports or trade less for the same imports while staying on a higher indifference curve.
\end{itemize}

\[
\boxed{\text{As an electronics exporter, a rise in }P_E/P_A \text{ raises welfare.}}
\]

\item \textbf{Export-biased growth and terms of trade.}

Export-biased growth: expansion of the PPF biased toward the export good (here $E$). At given world prices:
\begin{itemize}
\item Relative supply of $E$ rises.
\item For a \emph{large} country, this tends to \emph{lower} the world relative price of $E$.
\item This is a deterioration in the country's terms of trade (since it receives fewer imports per unit of $E$ exported).
\end{itemize}

The net welfare effect combines:
\begin{itemize}
\item A positive effect from larger production possibilities.
\item A negative terms-of-trade effect from lower export prices.
\end{itemize}

\[
\boxed{\text{Export-biased growth tends to worsen the exporter's terms of trade.}}
\]

\end{enumerate}

% ============================================================
\section*{Problem 5 — External Economies of Scale and Location}

\textbf{Setup:} Two countries A and B can produce semiconductors.  
External economies: average cost falls with industry size.  
A currently hosts the industry; B has identical technology but no initial industry.

\begin{enumerate}[label=(\arabic*)]

\item \textbf{Why can A remain the global producer?}

With external economies:
\begin{itemize}
\item Average (and marginal) cost declines with \emph{industry} output in a given country.
\item If A already has a large industry, its firms enjoy low costs.
\item If B has no industry, its initial costs are high (small scale).
\end{itemize}
Thus, even though technologies are identical, A's incumbency gives it a cost advantage due purely to scale and learning effects. New entrants in B cannot match A's costs at small scale, so production remains concentrated in A; this is \textbf{path dependence}.

\item \textbf{How can B enter despite A’s advantage?}

B can enter if:
\begin{itemize}
\item There is a sufficiently large \textbf{temporary demand shift} (e.g., local content requirement, a guaranteed market), allowing B's industry to start at a sufficiently large scale.
\item Or B provides \textbf{temporary support} (e.g., subsidies, public procurement) so that firms in B can operate at low effective costs while the industry scales up.
\item Once B's industry reaches critical size, its average costs fall, potentially matching or beating A's, allowing sustainable production even after support ends.
\end{itemize}

\item \textbf{Subsidy in B: shift in global production and welfare.}

If B subsidizes semiconductor production:
\begin{itemize}
\item \textbf{Short run:}  
  B's firms effectively face a lower cost curve. At world prices, they can profitably produce and export more, attracting more of the industry.
\item \textbf{Medium to long run:}  
  As B’s output expands, external economies lower its underlying average cost curve; A’s industry may shrink or disappear. Production can relocate fully to B.
\end{itemize}

\textbf{Welfare implications:}
\begin{itemize}
\item \textbf{For B:}  
  May gain in the long run if it captures a high-tech industry with strong external economies (higher wages, knowledge spillovers). The cost is the subsidy (fiscal burden) and possible short-run inefficiencies.
\item \textbf{For A:}  
  Loses an industry with external economies; wages and rents in that sector fall; may suffer lasting welfare loss.
\item \textbf{Global:}  
  If the new equilibrium has lower global average costs (e.g., more efficient location, larger eventual industry), world welfare can rise. But if relocation is driven mainly by policy rather than fundamentals, global welfare may not improve or may even fall (due to subsidy costs and trade distortions).
\end{itemize}

\textbf{Strategic industrial policy can shift location of an industry with external economies, with ambiguous global welfare effects.}


\end{enumerate}
\newpage
% ============================================================
\section*{Problem 6 — Trade Policy (Tariffs) for a Large Country}

\textbf{Domestic market (Home):}
\[
Q_D = 800 - 2P, \quad Q_S = P - 40.
\]
\textbf{Foreign export supply:}
\[
P^* = 150 + 0.5Q^*,
\]
where $P^*$ is the price received by foreign exporters; $Q^*$ is their export volume.  
$P$ is the domestic price in Home. A tariff $t$ drives a wedge:
\[
P = P^* + t.
\]

\begin{enumerate}[label=(\arabic*)]

\item \textbf{Free-trade price and imports.}

Under free trade: $t=0$, so $P = P^*$ and imports $Q = Q^*$.  

\textbf{Step 1: Derive Home’s import demand.}

Net imports:
\[
Q_M^D = Q_D - Q_S = (800 - 2P) - (P - 40) = 840 - 3P.
\]

\textbf{Step 2: Equilibrium with foreign export supply.}

Set import demand equal to export supply:
\[
Q = Q_M^D = 840 - 3P, \qquad P = P^* = 150 + 0.5Q.
\]

Substitute $Q=840 - 3P$ into export supply:
\[
P = 150 + 0.5(840 - 3P) = 150 + 420 - 1.5P = 570 - 1.5P.
\]
Solve:
\[
P + 1.5P = 570 \Rightarrow 2.5P = 570 \Rightarrow P^{FT} = 228.
\]

Free-trade imports:
\[
Q^{FT} = 840 - 3P^{FT} = 840 - 3\cdot 228 = 840 - 684 = 156.
\]

Also, domestic quantities:
\[
Q_D^{FT} = 800 - 2\cdot 228 = 344,\quad Q_S^{FT} = 228 - 40 = 188,
\]
and $Q_D^{FT} - Q_S^{FT} = 344 - 188 = 156$.

\[
\boxed{P^{FT} = 228,\quad Q^{FT} = 156.}
\]

\item \textbf{Tariff $t = 30$: new world price, domestic price, and imports.}

Now $P = P^* + t = P^* + 30$.

\textbf{Step 1: Express import demand in terms of $P^*$.}
\[
Q = Q_M^D = 840 - 3P = 840 - 3(P^* + 30) = 840 - 3P^* - 90 = 750 - 3P^*.
\]

\textbf{Step 2: Apply foreign export supply.}
\[
P^* = 150 + 0.5Q = 150 + 0.5(750 - 3P^*) = 150 + 375 - 1.5P^* = 525 - 1.5P^*.
\]
Solve:
\[
P^* + 1.5P^* = 525 \Rightarrow 2.5P^* = 525 \Rightarrow P^{*T} = 210.
\]

Then domestic price:
\[
P^{T} = P^{*T} + t = 210 + 30 = 240.
\]

Imports:
\[
Q^{T} = 750 - 3P^{*T} = 750 - 3\cdot 210 = 750 - 630 = 120.
\]

Check with $Q_D$ and $Q_S$:
\[
Q_D^{T} = 800 - 2\cdot 240 = 320,\quad Q_S^{T} = 240 - 40 = 200,\quad
Q_D^{T} - Q_S^{T} = 320 - 200 = 120.
\]

\[
\boxed{P^{*T} = 210,\quad P^{T} = 240,\quad Q^{T} = 120.}
\]

\item \textbf{Tariff revenue, terms-of-trade gain, deadweight loss.}

\paragraph{(a) Tariff revenue.}
\[
\text{Revenue} = t \cdot Q^{T} = 30 \times 120 = 3{,}600.
\]
\[
\boxed{\text{Tariff revenue} = \$3{,}600.}
\]

\paragraph{(b) Terms-of-trade (ToT) gain.}

Home’s terms of trade improve when the foreign price $P^*$ falls.  
Before tariff: $P^{FT} = 228$ (world price).  
After tariff: $P^{*T} = 210$. So:
\[
\Delta P_{\text{ToT}} = P^{FT} - P^{*T} = 228 - 210 = 18.
\]

Home still imports $Q^{T} = 120$ units at this lower foreign price. The ToT gain is approximately the rectangle:
\[
\text{ToT gain} = \Delta P_{\text{ToT}} \cdot Q^{T} = 18 \times 120 = 2{,}160.
\]

\[
\boxed{\text{Terms-of-trade gain} \approx \$2{,}160.}
\]

\paragraph{(c) Deadweight loss (DWL).}

We compute Home’s deadweight loss via changes in domestic quantities.

\textbf{Step 1: Changes in domestic quantities.}
\[
Q_D^{FT} = 344,\quad Q_D^{T} = 320 \Rightarrow \Delta Q_D = 344 - 320 = 24.
\]
\[
Q_S^{FT} = 188,\quad Q_S^{T} = 200 \Rightarrow \Delta Q_S = 200 - 188 = 12.
\]

\textbf{Step 2: Production distortion triangle.}
Base = $\Delta Q_S = 12$, height = tariff wedge $t = 30$:
\[
\text{DWL}_{\text{prod}} = \frac{1}{2} \cdot t \cdot \Delta Q_S
= \frac{1}{2} \cdot 30 \cdot 12 = 180.
\]

\textbf{Step 3: Consumption distortion triangle.}
Base = $\Delta Q_D = 24$, height = $t=30$:
\[
\text{DWL}_{\text{cons}} = \frac{1}{2} \cdot t \cdot \Delta Q_D
= \frac{1}{2} \cdot 30 \cdot 24 = 360.
\]

\textbf{Total DWL:}
\[
\text{DWL}_{\text{total}} = 180 + 360 = 540.
\]

\[
\boxed{\text{Total deadweight loss} = \$540.}
\]

\item \textbf{Does national welfare rise or fall?}

Net welfare effect for a large country:
\[
\Delta W = \text{ToT gain} - \text{DWL}.
\]

We found:
\[
\text{ToT gain} = 2{,}160,\quad \text{DWL} = 540.
\]

So:
\[
\Delta W = 2{,}160 - 540 = 1{,}620 > 0.
\]

\[
\boxed{\text{Home's national welfare rises by about \$1{,}620.}}
\]

\textbf{Economic intuition:}
\begin{itemize}
\item The tariff reduces imports, lowering the foreign export price from 228 to 210, improving Home’s terms of trade.
\item The gains from buying imports more cheaply (in world markets) exceed the domestic distortion losses from overproduction and underconsumption.
\item This is the classic ``optimal tariff'' logic for a large country.  
  In practice, retaliation by trading partners can undo such gains.
\end{itemize}

\end{enumerate}

\end{document}
